\documentclass[11pt, a4paper]{report}

\usepackage[utf8]{inputenc}
\usepackage[T1]{fontenc}
\usepackage[french]{babel}

% Appel des packages dans le dossier preambule
\def\packagepath{../../../../../preambule}   % path du package princial

\usepackage{\packagepath/page_layout} % <--- Nouveau
\usepackage{\packagepath/config_info}
\usepackage{\packagepath/cadres_cours}
\usepackage{\packagepath/zones_reponse}
\usepackage{\packagepath/config_ds}
\usepackage{\packagepath/config_maths}

% --- PILOTAGE ---
%\def\visibleornot{visible}    % visible or invisible
\ifdefined\visibleornot\else\def\visibleornot{visible}\fi


\def\documentpath{./Sources_Latex}
\graphicspath{{./Images/}}


\def\ptsexoA{0}



\begin{document}

\def\level{Premiere}  % Classe à droite
\def\course{NSI}
\def\chaptername{Les listes et tuples} % Titre au centre
\def\docname{AP05~--~TD 1}        % Identifiant à gauche

\pagestyle{DS_LP}

\NomPrenom{}

\begin{exercice_cours}[pour chaque instruction, indiquer ce qui est affiché: \quad \texttt{L = [2, 6, 4, 8, 10]}]{}
    
\begin{minipage}{0.33\linewidth}
    \begin{tabular}{p{2cm}c}
    \hline
    Instruction & Affichage \\    \hline
    \texttt{L[0]} & \reponseLigne[2.cm]{\texttt{2}}  \\
    \texttt{L[2]} & \reponseLigne[2.cm]{\texttt{4}}  \\
    \hline  
    \end{tabular}
\end{minipage}\hfill{}
\begin{minipage}{0.33\linewidth}
\begin{tabular}{p{2cm}c}
    \hline
    Instruction & Affichage \\    \hline
    \texttt{L[-1]} & \reponseLigne[2.cm]{\texttt{10}}  \\
    \texttt{L[1:4]} & \reponseLigne[2.cm]{\texttt{[6, 4, 8]}}  \\
    \hline  
    \end{tabular}

\end{minipage}\hfill{}
\begin{minipage}{0.33\linewidth}
\begin{tabular}{p{2cm}c}
    \hline
    Instruction & Affichage \\    \hline
    \texttt{L[:3]} & \reponseLigne[2.cm]{\texttt{[2, 6, 4]}}  \\
    \texttt{L[2:]} & \reponseLigne[2.cm]{\texttt{[4, 8, 10]}}  \\
    \hline  
    \end{tabular}
    
\end{minipage}     
\end{exercice_cours}

\begin{exercice_cours}[pour chaque instruction, indiquer ce qui est affiché: \quad \texttt{M = [[2, 6, 4], [8, 10]]}]{}
    
\begin{minipage}{0.33\linewidth}
    \begin{tabular}{p{2cm}c}
    \hline
    Instruction & Affichage \\    \hline
    \texttt{M[0]} & \reponseLigne[2.cm]{\texttt{[2, 6, 4]}}  \\
    \texttt{M[1]} & \reponseLigne[2.cm]{\texttt{[8, 10]}}  \\
    \hline  
    \end{tabular}
\end{minipage}\hfill{}
\begin{minipage}{0.33\linewidth}
\begin{tabular}{p{2cm}c}
    \hline
    Instruction & Affichage \\    \hline
    \texttt{M[1][0]} & \reponseLigne[2.cm]{\texttt{8}}  \\
    \texttt{M[0][1:3]} & \reponseLigne[2.cm]{\texttt{[6, 4]}}  \\
    \hline  
    \end{tabular}

\end{minipage}\hfill{}
\begin{minipage}{0.33\linewidth}
\begin{tabular}{p{2cm}c}
    \hline
    Instruction & Affichage \\    \hline
    \texttt{len(M)} & \reponseLigne[2.cm]{\texttt{2}}  \\
    \texttt{len(M[0])} & \reponseLigne[2.cm]{\texttt{3}}  \\
    \hline  
    \end{tabular}
    
\end{minipage}     
\end{exercice_cours}


\begin{minipage}{0.48\linewidth}
        \begin{exercice_cours}[Ecrire un code]
        Ecrire une fonction \verb|multiples(n)| qui renvoie la liste des multiples de 3 inférieurs à n.
    \end{exercice_cours}

    \end{minipage}\hfill
    \begin{minipage}{0.48\linewidth}
\begin{blocvisible}[height=4]
    \begin{CodePythontexAlt}[TaillePolice=\large, Largeur=1\linewidth,Centre,Lignes=true,PremLigne=1]{}
def multiples(n):
    result = []
    for i in range(n):
        if i % 3 == 0:
            result.append(i)
    return result
    \end{CodePythontexAlt}
\end{blocvisible}
    \end{minipage}



    \begin{minipage}{0.48\linewidth}
        \begin{exercice_cours}[Ecrire un code]
        Ecrire une fonction \verb|remplace_pair(liste)| qui prend en paramètre une liste de nombres et qui remplace tous les éléments pairs de la liste par 0. La liste doit être modifiée en place.
    \end{exercice_cours}

    \end{minipage}\hfill
    \begin{minipage}{0.48\linewidth}
\begin{blocvisible}[height=3.]
    \begin{CodePythontexAlt}[TaillePolice=\large, Largeur=1\linewidth,Centre,Lignes=true,PremLigne=1]{}
def remplace_pair(liste):
    for i in range(len(liste)):
        if liste[i] % 2 == 0:
            liste[i] = 0
    \end{CodePythontexAlt}
\end{blocvisible}
    \end{minipage}

    
    \begin{minipage}{0.48\linewidth}
        \begin{exercice_cours}[Test d'appartenance]

Ecrire la fonction  qui prend en paramètre une liste et une valeur, et qui retourne: 
\begin{itemize}
        \item \texttt{True} si la valeur est présente dans la liste
        \item \texttt{False} sinon 
    \end{itemize}

    Fonction: \verb|est_present(liste, valeur)|
    \end{exercice_cours}

    \end{minipage}\hfill
    \begin{minipage}{0.48\linewidth}
\begin{blocvisible}[height=4.]
    \begin{CodePythontexAlt}[TaillePolice=\large, Largeur=1\linewidth,Centre,Lignes=true,PremLigne=1]{}
def est_present(liste, valeur):
    for element in liste:
        if element == valeur:
            return True
    return False
    \end{CodePythontexAlt}
\end{blocvisible}


    \end{minipage}

        
\begin{minipage}{0.45\linewidth}
           \begin{exercice_cours}[Séries de valeurs]

Ecrire la fonction \texttt{serie(liste)} qui prend en paramètre une liste et qui retourne le nombre de \texttt{1} consécutifs à gauche de la liste. Par exemple, pour la liste \texttt{[1, 1, 0, 1, 1, 1, 1, 1, 1]}, la fonction doit retourner \texttt{2}. 
    \end{exercice_cours}

    \end{minipage}\hfill
    \begin{minipage}{0.53\linewidth}
\begin{blocvisible}[height=5]
    \begin{CodePythontexAlt}[TaillePolice=\large, Largeur=1\linewidth,Centre,Lignes=true,PremLigne=1]{}
def serie(liste):
    count = 0
    for element in liste:
        if element == 1:
            count += 1
        else:
            return count
    return count
    
    \end{CodePythontexAlt}
\end{blocvisible}

    \end{minipage}

\end{document}